\par\Needspace{8\baselineskip}\bigskip
\pdfbookmark[3]{Exposition, 16}{pallas-II-article16}
\begin{center}16.\end{center}
\noindent\hspace*{\parindent}If instead of \(T,V\) we prefer to use \(R,S\), the most convenient form of the calculation seems to be to refer the positions of both the perturbing and perturbed planets to three mutually perpendicular planes. One is the plane of the osculating orbit of the perturbed planet, and one of the other two passes through its line of apsides. Let \(X',Y',Z'\) be the coordinates of the perturbing planet and \(X,Y,0\) those of the perturbed planet, given by
\SourceLeaf{466}
\[
\begin{aligned}
X'&=r'\cos PL',&\qquad X&=r\cos(v-\tilde\omega)\\
Y'&=r'\cos\Pi L',&Y&=r\sin(v-\tilde\omega)\\
Z'&=r'\cos ML',&&
\end{aligned}
\]
and we have
\[
\begin{gathered}
\rho=\BareRoot{((X'-X)^2+(Y'-Y)^2+Z'Z')}\\[4pt]
\begin{aligned}
\dfrac{R}{\mu}&=\dfrac{X'-X}{\rho^3}-\dfrac{X'}{r'^2}\\
\dfrac{S}{\mu}&=\dfrac{Y'-Y}{\rho^3}-\dfrac{Y'}{r'^3}\\
\dfrac{W}{\mu}&=\dfrac{Z'}{\rho^3}-\dfrac{Z'}{r'^3}\mathpunct{\cdot}
\end{aligned}
\end{gathered}
\]
\par The most convenient way of computing these coordinates is sufficiently well known that there is no need to dwell on it here.
\par In all the preceding investigations, the various arcs are understood to be expressed in units of the radius. In all homogeneous formulas, where the arcs have the same dimension on both sides, they may therefore also be expressed in arcseconds. In formulas that are inhomogeneous, the factor 206265 must be introduced on the side where the arcs have one dimension less. Thus equations 61, 62, 66, 68, 70, 74, 75, 76 remain the same, whereas in equations 63, 64, 65, 72, 73 the right-hand side must be divided by 206265 if all the arcs are understood to be expressed in arcseconds. If, however, one were to exempt the arc \(n\) and continue to express it in units of the radius, the right-hand side of equations 61, 62, 66, etc. would have to be multiplied by 206265, while the other equations 63, 64, 65, etc. would remain unchanged.
\par\medskip\begin{center}\rule{4em}{.4pt}\end{center}
